% Please add the following required packages to your document preamble:
% \usepackage{multirow}
% \usepackage{graphicx}
\begin{table}[t!]
\centering
\begin{tabular}{l|l|l}
\hline
\textbf{Backbone} & \textbf{Method}     & \textbf{Acc. (\%)} \\ \hline
gpt-3.5     & Zero-shot & 31.2               \\ \hline
\multirow{2}{*}{\begin{tabular}[c]{@{}c@{}}llama3\end{tabular}} & Zero-shot & 33.0 \\
                  & Fine-tuning        & 36.5               \\ \hline
\multirow{2}{*}{\begin{tabular}[c]{@{}c@{}}gemma \end{tabular}}     & Zero-shot & 36.9 \\
                  & Fine-tuning        & 46.7               \\ \hline
\end{tabular}%
\caption{Performance of answerability assessment models on RaCGEval benchmarks.}
\label{tab:2}
\end{table}

\section{Experiments}
\subsection{Experimental settings}
In this experiment, we use the following three LLMs: \textit{gemma-1.1-7b-it} \cite{team2024gemma}, \textit{llama3-instruct-8b}~\cite{dubey2024llama3herdmodels}, and \textit{gpt-3.5-turbo-0613}. For fine-tuned models, we trained them to predict one of the following three tokens for a given prompt: "A (answerable)", "U (unanswerable)" and "P (partially answerable)". Therefore, during inference, we compare the likelihoods of the three tokens to make a prediction. We employ the quantized low-Rank adaptation (QLoRA)~\cite{tim2024qlora} for parameter-efficient fine-tuning to reduce training VRAM.

\subsection{Experimental results}
We measure the performance of the answerability assessment models on the RaCGEval benchmark dataset. Table~\ref{tab:2} shows the performance of zero-shot inference or fine-tuning across various LLMs.

For zero-shot inference, the accuracy is similar to the chance level of 3-way multi-class classification across all LLM backbones. This indicates that the answerability assessment task in the code domain is not familiar to several LLMs. However, it is evident that substantial performance improvements can be achieved through fine-tuning. Nevertheless, there is still room for improvement, and we will discuss potential solutions in the discussion session.